%% informe_practica1.tex
%% Práctica 1 — Jerarquía de Memoria y Comportamiento de Caché
%% Arquitectura de Computadores — FinisTerrae III (CESGA)
%% Compilar: pdflatex informe_practica1.tex  (2 pasadas)

\documentclass[12pt,a4paper]{article}

% ── Paquetes ────────────────────────────────────────────────────────────────
\usepackage[utf8]{inputenc}
\usepackage[T1]{fontenc}
\usepackage{geometry}
\geometry{left=2.5cm, right=2.5cm, top=2.8cm, bottom=2.8cm, headheight=14pt}

\usepackage{amsmath}
\usepackage{graphicx}
\usepackage{booktabs}
\usepackage{multirow}
\usepackage{array}
\usepackage{xcolor}
\usepackage{listings}
\usepackage{caption}
\usepackage{subcaption}
\usepackage{hyperref}
\usepackage{fancyhdr}
\usepackage{titlesec}
\usepackage{enumitem}
\usepackage{float}
\usepackage{parskip}

% ── Colores personalizados ──────────────────────────────────────────────────
\definecolor{verde}{RGB}{40,167,69}
\definecolor{amarillo}{RGB}{255,193,7}
\definecolor{rojo}{RGB}{220,53,69}
\definecolor{naranja}{RGB}{253,126,20}
\definecolor{azul}{RGB}{0,86,179}
\definecolor{grisclaro}{RGB}{248,249,250}
\definecolor{griscode}{RGB}{40,40,40}

% ── Configuración de listings (código C) ────────────────────────────────────
\lstset{
  language=C,
  basicstyle=\ttfamily\small,
  keywordstyle=\color{azul}\bfseries,
  commentstyle=\color{verde}\itshape,
  stringstyle=\color{rojo},
  numbers=left,
  numberstyle=\tiny\color{gray},
  numbersep=8pt,
  frame=single,
  backgroundcolor=\color{grisclaro},
  breaklines=true,
  captionpos=b,
  tabsize=4,
  showstringspaces=false,
  xleftmargin=15pt,
}

% ── Cabecera y pie de página ─────────────────────────────────────────────────
\pagestyle{fancy}
\fancyhf{}
\fancyhead[L]{\small Arquitectura de Computadores — Práctica 1}
\fancyhead[R]{\small FinisTerrae III · Intel Xeon Platinum 8352Y}
\fancyfoot[C]{\thepage}
\renewcommand{\headrulewidth}{0.4pt}

% ── Formato de secciones ────────────────────────────────────────────────────
\titleformat{\section}{\large\bfseries\color{azul}}{{\thesection}}{1em}{}[\vspace{-4pt}\rule{\linewidth}{0.6pt}]
\titleformat{\subsection}{\normalsize\bfseries}{{\thesubsection}}{1em}{}

% ── Macros de conveniencia ──────────────────────────────────────────────────
\newcommand{\ciclos}[1]{\textbf{#1} ciclos}
\newcommand{\verde}[1]{\textcolor{verde}{#1}}
\newcommand{\amarillo}[1]{\textcolor{orange}{#1}}
\newcommand{\rojo}[1]{\textcolor{rojo}{#1}}

% ── Metadatos PDF ───────────────────────────────────────────────────────────
\hypersetup{
  pdftitle={Práctica 1 — Jerarquía de Memoria},
  pdfauthor={Arquitectura de Computadores},
  colorlinks=true,
  linkcolor=azul,
  urlcolor=azul,
}

% ════════════════════════════════════════════════════════════════════════════
\begin{document}

% ── Portada ─────────────────────────────────────────────────────────────────
\begin{titlepage}
  \centering
  \vspace*{2cm}
  {\LARGE\bfseries\color{azul} Arquitectura de Computadores\par}
  \vspace{0.4cm}
  {\large Práctica 1\par}
  \vspace{1.2cm}
  \rule{0.8\textwidth}{1pt}
  \vspace{0.5cm}

  {\LARGE\bfseries Jerarquía de Memoria y Comportamiento de Caché\par}
  \vspace{0.3cm}
  {\large Estudio del efecto de la localidad de los accesos a memoria\par
   en las prestaciones de programas en microprocesadores\par}

  \vspace{0.5cm}
  \rule{0.8\textwidth}{1pt}
  \vspace{2cm}

  \begin{tabular}{ll}
    \textbf{Plataforma:} & FinisTerrae III (CESGA) \\[4pt]
    \textbf{Procesador:} & Intel Xeon Platinum 8352Y (Ice Lake) \\[4pt]
    \textbf{Compilador:} & \texttt{gcc -O0} \\[4pt]
    \textbf{Jobs Slurm:}  & 5366644, 5367628, 5367629, 5367630 \\[4pt]
    \textbf{Métrica:}     & Media geométrica de los 3 mejores de 10 repeticiones \\
  \end{tabular}

  \vfill
  {\large\today\par}
\end{titlepage}

% ── Índice ───────────────────────────────────────────────────────────────────
\tableofcontents
\clearpage

% ════════════════════════════════════════════════════════════════════════════
\section{Introducción y objetivos}

Esta práctica estudia el efecto de la jerarquía de memoria sobre el coste temporal de los accesos a datos. El objetivo principal es \textbf{caracterizar experimentalmente la latencia de cada nivel de caché} del procesador Intel Xeon Platinum 8352Y variando dos parámetros clave:

\begin{itemize}[leftmargin=1.5em]
  \item \textbf{D (stride):} separación en elementos entre posiciones accedidas. Controla la localidad espacial.
  \item \textbf{L (líneas):} número de líneas de caché distintas referenciadas. Determina el nivel de la jerarquía que almacena los datos.
\end{itemize}

El experimento base suma $R$ elementos del vector \texttt{A[]} con patrón $A[0], A[D], A[2D], \ldots, A[(R{-}1)D]$, accedidos de forma \textbf{indirecta} a través de un vector de índices \texttt{ind[]}. Se añaden dos experimentos adicionales: acceso con tipo \texttt{int} en lugar de \texttt{double}, y acceso \textbf{directo} eliminando el vector de índices.

% ════════════════════════════════════════════════════════════════════════════
\section{Arquitectura del sistema}

\subsection{Jerarquía de memoria del Intel Xeon Platinum 8352Y}

El procesador implementa la microarquitectura \textbf{Sunny Cove (Ice Lake-SP)} con la siguiente jerarquía:

\begin{table}[H]
\centering
\caption{Jerarquía de caché del Intel Xeon Platinum 8352Y}
\label{tab:jerarquia}
\begin{tabular}{lllll}
\toprule
\textbf{Nivel} & \textbf{Capacidad} & \textbf{Asociatividad} & \textbf{S (líneas)} & \textbf{Latencia aprox.} \\
\midrule
L1d & 48 KB & 12-way & $S_1 = 768$ & $\sim$5 ciclos \\
L2  & 1.25 MB & 20-way & $S_2 = 20\,480$ & $\sim$14 ciclos \\
L3  & 48 MB & 12-way & $S_3 = 786\,432$ & $\sim$40--50 ciclos \\
RAM & 256 GB DDR4 & — & — & $\sim$150--200 ciclos \\
\bottomrule
\end{tabular}
\end{table}

La línea de caché tiene un tamaño de \textbf{64 bytes} en todos los niveles. Los experimentos se ejecutan en \textbf{un único core} para garantizar que L1 y L2 son privadas y no hay contención con otros procesos.

\subsection{Prefetcher hardware de Ice Lake}

El procesador dispone de \textbf{cuatro prefetchers hardware} independientes que actúan coordinadamente:

\subsubsection*{Prefetchers de L1 (cargan desde L2 hacia L1)}

\begin{description}[leftmargin=1.5em]
  \item[DCU Streamer (Next-Line Prefetcher):] Detecta accesos ascendentes y precarga automáticamente la siguiente línea de caché. Se activa con cualquier patrón secuencial.
  \item[DCU IP-based Stride Prefetcher:] Rastrea instrucciones de carga individuales. Cuando detecta un stride regular, emite precargas a \texttt{dirección\_actual + stride}. Detecta strides de hasta 2 KB en ambas direcciones.
\end{description}

\subsubsection*{Prefetchers de L2 (cargan desde L3 hacia L2)}

\begin{description}[leftmargin=1.5em]
  \item[L2 Streamer:] Monitoriza secuencias de peticiones desde L1. Puede correr hasta \textbf{20 líneas por delante} y gestiona hasta 32 streams simultáneos. No cruza límites de página de 4 KB.
  \item[L2 Spatial (Adjacent Cache Line) Prefetcher:] Completa cada línea cargada en L2 con su línea adyacente, garantizando bloques de 128 bytes alineados.
\end{description}

\subsubsection*{Implicaciones experimentales}

Con $D=1$, los cuatro prefetchers actúan coordinadamente ocultando completamente las latencias de L2, L3 e incluso RAM. Con $D=16$ se produce \emph{prefetch pollution}: el prefetcher spatial trae líneas adyacentes inútiles duplicando el tráfico efectivo. Con $D \geq 64$, el stride supera el umbral eficiente del L2 Streamer, aunque el IP Stride Prefetcher sigue actuando parcialmente.

% ════════════════════════════════════════════════════════════════════════════
\section{Metodología}

\subsection{Parámetros experimentales}

\subsubsection*{Valores de L}

Se toman 11 medidas abarcando toda la jerarquía, desde L1 hasta RAM profunda:

\begin{table}[H]
\centering
\caption{Valores de L utilizados y nivel de jerarquía correspondiente}
\label{tab:valoresL}
\begin{tabular}{rrlr}
\toprule
\textbf{L (líneas)} & \textbf{Fracción} & \textbf{Footprint} & \textbf{Nivel} \\
\midrule
384       & $0.5 \times S_1$    & 24 KB   & L1 (holgado) \\
1\,152    & $1.5 \times S_1$    & 72 KB   & L2 (desborda L1) \\
10\,240   & $0.5 \times S_2$    & 640 KB  & L2 (holgado) \\
15\,360   & $0.75 \times S_2$   & 960 KB  & L2 (ajustado) \\
40\,960   & $2 \times S_2$      & 2.5 MB  & L3 (desborda L2) \\
81\,920   & $4 \times S_2$      & 5 MB    & L3 \\
163\,840  & $8 \times S_2$      & 10 MB   & L3 \\
524\,288  & $0.67 \times S_3$   & 32 MB   & L3 (llenando) \\
786\,432  & $1 \times S_3$      & 48 MB   & Límite exacto L3 \\
1\,572\,864 & $2 \times S_3$   & 96 MB   & \textbf{RAM real} \\
3\,145\,728 & $4 \times S_3$   & 192 MB  & \textbf{RAM profunda} \\
\bottomrule
\end{tabular}
\end{table}

\subsubsection*{Valores de D (stride)}

Se emplean 5 valores potencia de 2: $D \in \{1, 8, 16, 64, 128\}$, seleccionados para cubrir todos los regímenes del prefetcher:

\begin{itemize}[leftmargin=1.5em]
  \item $D=1$: acceso completamente secuencial; localidad espacial máxima.
  \item $D=8$: stride de 64 bytes; exactamente una línea de caché por acceso.
  \item $D=16$: stride de 128 bytes; activa el prefetch pollution del spatial prefetcher.
  \item $D=64, D=128$: strides grandes que superan el umbral del L2 Streamer.
\end{itemize}

\subsection{Cálculo de R}

Para un valor de D y L dados, el número de elementos a sumar $R$ se calcula como:
\[
  R = \frac{L \times \text{CLS}}{D \times \text{sizeof(tipo)}}
\]
donde CLS = 64 bytes. El vector \texttt{A[]} tiene $N = (R-1) \cdot D + 1$ elementos en total.

\subsection{Tratamiento estadístico}

Para cada combinación $(D, L)$ se ejecutan \textbf{10 repeticiones externas}, cada una con 10 repeticiones internas del bucle de suma. Se seleccionan los 3 menores valores de ciclos/acceso y se calcula su \textbf{media geométrica}, eliminando medidas atípicas por interferencia del SO.

\subsection{Variantes del experimento}

\begin{enumerate}[leftmargin=1.5em]
  \item \textbf{Double + indirecto} (\texttt{acp1.c}): caso base con \texttt{double} y acceso \texttt{A[ind[i]]}.
  \item \textbf{Int + indirecto} (\texttt{acp1\_int.c}): tipo \texttt{int} (4 bytes); R es el doble para el mismo L.
  \item \textbf{Double + directo} (\texttt{acp1\_directo.c}): elimina \texttt{ind[]} y accede como \texttt{A[i*D]}.
\end{enumerate}

En total se realizan $3 \times 5 \times 11 \times 10 = \mathbf{1\,650}$ mediciones.

% ════════════════════════════════════════════════════════════════════════════
\section{Implementación}

\subsection{Programa principal: \texttt{acp1.c}}

El programa recibe como argumentos el stride D y el número de líneas L. Los aspectos más relevantes de la implementación son:

\begin{itemize}[leftmargin=1.5em]
  \item \textbf{Alineación de memoria:} \texttt{aligned\_alloc(64, N * sizeof(double))} garantiza que el inicio de \texttt{A[]} coincide con el inicio de una línea de caché. Esto es crítico para resultados precisos en la zona L1, donde la alineación determina si el primer acceso produce un fallo de caché parcial.
  \item \textbf{Calentamiento implícito:} la inicialización de \texttt{A[]} recorre las mismas posiciones que el bucle de medición, pre-cargando las líneas relevantes en caché y evitando efectos de arranque en frío de las TLBs.
  \item \textbf{Evitar optimizaciones:} el resultado de la suma se almacena en \texttt{acum[k]} y se usa en una condición al final del programa, impidiendo que \texttt{gcc -O0} elimine el bucle por dead-code elimination.
  \item \textbf{Medición:} se miden los ciclos totales de las 10 repeticiones internas con las rutinas \texttt{start\_counter()} / \texttt{get\_counter()} de \texttt{counter.h}, y se divide por $R \times 10$ para obtener ciclos por acceso.
\end{itemize}

\begin{lstlisting}[caption={Fragmento principal de \texttt{acp1.c}: cálculo de R y bucle de medición}, label=lst:acp1]
// Cálculo de R: líneas referenciadas × bytes/línea / (stride × bytes/elemento)
long long R = (long long)L * CLS / (D * sizeof(double));
long long N = (R - 1) * D + 1;

// Reserva alineada al inicio de línea de caché
double *A = (double*)aligned_alloc(CLS, N * sizeof(double));
int    *ind = (int*)malloc(R * sizeof(int));

// Inicialización: calentamiento implícito de caché y TLBs
for (long long i = 0; i < R; i++) ind[i] = i * D;
for (long long i = 0; i < R; i++) A[ind[i]] = (double)rand() / RAND_MAX;

// Sección crítica de medición
start_counter();
for (int k = 0; k < REPS; k++) {
    double suma = 0.0;
    for (long long i = 0; i < R; i++)
        suma += A[ind[i]];      // Acceso indirecto
    acum[k] = suma;
}
double ciclos_totales = get_counter();

double ciclos_por_acceso = ciclos_totales / ((double)R * REPS);
\end{lstlisting}

% ════════════════════════════════════════════════════════════════════════════
\section{Resultados experimentales}

\subsection{Experimento base: Double + acceso indirecto}

\begin{table}[H]
\centering
\caption{Ciclos de CPU por acceso — Double + acceso indirecto (\texttt{acp1.c}). Sombreado: \verde{$<8$~cic.} / \amarillo{8--12~cic.} / \rojo{$>12$~cic.}}
\label{tab:double_ind}
\setlength{\tabcolsep}{5pt}
\begin{tabular}{rllccccc}
\toprule
\textbf{L} & \textbf{Nivel} & \textbf{Footprint} & \textbf{D=1} & \textbf{D=8} & \textbf{D=16} & \textbf{D=64} & \textbf{D=128} \\
\midrule
384        & L1         & 24 KB   & \verde{7.08} & \verde{7.14} & \verde{7.06} & \verde{7.15} & \verde{7.12} \\
1\,152     & L2         & 72 KB   & \verde{7.09} & \verde{7.17} & \verde{7.18} & \verde{7.20} & \verde{7.28} \\
10\,240    & L2         & 640 KB  & \verde{7.12} & \verde{7.19} & \verde{7.21} & \verde{7.29} & \verde{7.27} \\
15\,360    & L2         & 960 KB  & \verde{7.13} & \verde{7.20} & \verde{7.23} & \verde{7.41} & \verde{7.33} \\
40\,960    & L3         & 2.5 MB  & \verde{7.15} & \verde{7.45} & \amarillo{8.96} & \amarillo{8.18} & \amarillo{8.28} \\
81\,920    & L3         & 5 MB    & \verde{7.18} & \verde{7.52} & \amarillo{9.48} & \amarillo{8.23} & \amarillo{8.21} \\
163\,840   & L3         & 10 MB   & \verde{7.22} & \amarillo{8.20} & \amarillo{10.32} & \amarillo{8.33} & \amarillo{8.64} \\
524\,288   & L3         & 32 MB   & \verde{7.22} & \amarillo{11.01} & \rojo{17.59} & \rojo{17.14} & \rojo{17.09} \\
786\,432   & Límite L3  & 48 MB   & \verde{7.23} & \amarillo{11.15} & \rojo{18.30} & \rojo{18.43} & \rojo{18.59} \\
1\,572\,864 & \textbf{RAM} & 96 MB & \verde{7.23} & \amarillo{11.09} & \rojo{18.99} & \rojo{20.55} & \rojo{19.42} \\
3\,145\,728 & \textbf{RAM} & 192 MB & \verde{7.22} & \amarillo{11.10} & \rojo{18.51} & \rojo{20.86} & \rojo{19.68} \\
\bottomrule
\end{tabular}
\end{table}

\subsection{Experimento adicional 1: Int + acceso indirecto}

\begin{table}[H]
\centering
\caption{Ciclos de CPU por acceso — Int + acceso indirecto (\texttt{acp1\_int.c})}
\label{tab:int_ind}
\setlength{\tabcolsep}{5pt}
\begin{tabular}{rllccccc}
\toprule
\textbf{L} & \textbf{Nivel} & \textbf{Footprint} & \textbf{D=1} & \textbf{D=8} & \textbf{D=16} & \textbf{D=64} & \textbf{D=128} \\
\midrule
384        & L1         & 24 KB   & \verde{6.77} & \verde{6.79} & \verde{6.84} & \verde{6.80} & \verde{6.82} \\
1\,152     & L2         & 72 KB   & \verde{6.78} & \verde{6.82} & \verde{7.17} & \verde{7.23} & \verde{7.21} \\
10\,240    & L2         & 640 KB  & \verde{6.79} & \verde{6.83} & \verde{7.18} & \verde{7.23} & \verde{7.28} \\
15\,360    & L2         & 960 KB  & \verde{6.80} & \verde{6.87} & \verde{7.24} & \verde{7.24} & \verde{7.22} \\
40\,960    & L3         & 2.5 MB  & \verde{6.80} & \verde{6.88} & \verde{7.26} & \amarillo{8.00} & \verde{7.94} \\
81\,920    & L3         & 5 MB    & \verde{6.80} & \verde{6.91} & \verde{7.28} & \verde{7.73} & \verde{7.76} \\
163\,840   & L3         & 10 MB   & \verde{6.80} & \verde{6.93} & \verde{7.32} & \verde{7.78} & \verde{7.76} \\
524\,288   & L3         & 32 MB   & \verde{6.81} & \verde{7.69} & \amarillo{8.58} & \rojo{15.88} & \rojo{14.92} \\
786\,432   & Límite L3  & 48 MB   & \verde{6.81} & \verde{7.79} & \amarillo{8.80} & \rojo{16.87} & \rojo{16.01} \\
1\,572\,864 & \textbf{RAM} & 96 MB & \verde{6.81} & \verde{7.87} & \amarillo{9.27} & \rojo{17.04} & \rojo{18.35} \\
3\,145\,728 & \textbf{RAM} & 192 MB & \verde{6.81} & \verde{7.89} & \amarillo{9.40} & \rojo{17.53} & \rojo{19.71} \\
\bottomrule
\end{tabular}
\end{table}

\subsection{Experimento adicional 2: Double + acceso directo}

\begin{table}[H]
\centering
\caption{Ciclos de CPU por acceso — Double + acceso directo (\texttt{acp1\_directo.c})}
\label{tab:double_dir}
\setlength{\tabcolsep}{5pt}
\begin{tabular}{rllccccc}
\toprule
\textbf{L} & \textbf{Nivel} & \textbf{Footprint} & \textbf{D=1} & \textbf{D=8} & \textbf{D=16} & \textbf{D=64} & \textbf{D=128} \\
\midrule
384        & L1         & 24 KB   & \verde{7.06} & \verde{6.98} & \verde{6.82} & \verde{6.42} & \verde{6.36} \\
1\,152     & L2         & 72 KB   & \verde{7.08} & \verde{7.09} & \verde{7.07} & \verde{6.88} & \verde{6.81} \\
10\,240    & L2         & 640 KB  & \verde{7.10} & \verde{7.15} & \verde{7.15} & \verde{7.10} & \verde{7.08} \\
15\,360    & L2         & 960 KB  & \verde{7.11} & \verde{7.18} & \verde{7.14} & \verde{7.12} & \verde{7.09} \\
40\,960    & L3         & 2.5 MB  & \verde{7.13} & \verde{7.35} & \verde{7.57} & \verde{7.50} & \verde{7.48} \\
81\,920    & L3         & 5 MB    & \verde{7.12} & \verde{7.30} & \verde{7.53} & \verde{7.55} & \verde{7.32} \\
163\,840   & L3         & 10 MB   & \verde{7.12} & \verde{7.32} & \verde{7.45} & \verde{7.56} & \verde{7.41} \\
524\,288   & L3         & 32 MB   & \verde{7.14} & \amarillo{8.36} & \rojo{12.63} & \rojo{13.69} & \rojo{12.99} \\
786\,432   & Límite L3  & 48 MB   & \verde{7.15} & \amarillo{8.54} & \rojo{14.62} & \rojo{15.13} & \rojo{13.91} \\
1\,572\,864 & \textbf{RAM} & 96 MB & \verde{7.16} & \amarillo{8.58} & \rojo{16.25} & \rojo{17.22} & \rojo{15.44} \\
3\,145\,728 & \textbf{RAM} & 192 MB & \verde{7.16} & \amarillo{8.56} & \rojo{16.64} & \rojo{18.72} & \rojo{16.59} \\
\bottomrule
\end{tabular}
\end{table}

% ════════════════════════════════════════════════════════════════════════════
\section{Análisis por zonas de la jerarquía}

\subsection{Zona L1 — $L \leq 768$ líneas (footprint $\leq$ 48 KB)}

En la zona L1, \textbf{todas las curvas convergen} en el rango 6.4--7.2 ciclos/acceso, independientemente del stride D y del tipo de dato. Con 24 KB de footprint, el conjunto de datos cabe completamente en la caché L1d de 48 KB y cada acceso se resuelve sin fallo de caché.

La ligera ventaja del acceso directo con strides grandes (D=64: 6.42 ciclos frente a 7.15 del indirecto) se debe a que sin \texttt{ind[]}, todo el espacio L1 se dedica exclusivamente a \texttt{A[]}. La ventaja del tipo \texttt{int} ($\sim$0.3 ciclos menos) refleja que la instrucción de suma entera \texttt{ADD} tiene menor latencia que \texttt{FADD} de punto flotante.

\textbf{Conclusión:} en la zona L1 el stride es irrelevante. Todos los accesos tienen coste $\approx 5$--7 ciclos (compatible con la latencia teórica de L1 más overhead del bucle sin optimización).

\subsection{Zona L2 — $768 < L \leq 20\,480$ líneas (footprint 72 KB -- 960 KB)}

Al superar L1, la latencia teórica de L2 es $\sim$14 ciclos. Sin embargo, los valores medidos permanecen por debajo de 7.5 ciclos en todas las variantes. Esto confirma la actuación del \textbf{DCU IP Stride Prefetcher}, que detecta el stride constante del bucle y emite precargas desde L2 hacia L1 antes de que sean necesarias, ocultando prácticamente toda la latencia.

La penalización al pasar de L1 a L2 es inferior a 0.2 ciclos, lo que evidencia la eficacia del prefetcher incluso para strides tan grandes como D=128.

\subsection{Zona L3 — $20\,480 < L \leq 786\,432$ líneas (footprint 2.5 MB -- 48 MB)}

Esta zona muestra el comportamiento más diferenciado entre strides:

\begin{description}[leftmargin=1.5em]
  \item[$D=1$:] Permanece completamente plano en $\sim$7.2 ciclos hasta el límite de L3. El \textbf{L2 Streamer} (20 líneas por delante) oculta la latencia de L3 de forma idéntica a la de L2.

  \item[$D=8$:] Crece progresivamente de 7.5 hasta $\sim$11 ciclos al llegar al límite de L3. El prefetcher sigue siendo eficaz pero empieza a limitarse porque el stride de 512 bytes le obliga a emitir precargas más espaciadas.

  \item[$D=16$:] Peor caso observado. Escala de 9 hasta $\sim$18 ciclos. El \textbf{L2 Spatial Prefetcher} trae líneas adyacentes (stride real = 1024 bytes), generando \emph{prefetch pollution}: tráfico inútil que satura el ancho de banda L2-L3 y perjudica los accesos reales.

  \item[$D=64, D=128$:] El stride supera ampliamente el umbral del L2 Streamer. Paradójicamente, alcanzan latencias similares a D=16 en la zona L3 profunda ($\sim$18 ciclos) porque ahora el prefetcher reduce su actividad, dejando que las latencias de L3 se manifiesten sin la polución de D=16.
\end{description}

\subsection{Zona RAM — $L > 786\,432$ líneas (footprint $>$ 48 MB)}

La transición L3$\to$RAM es el hallazgo más relevante de los experimentos ampliados. Entre $L=786\,432$ y $L=1\,572\,864$ se observa una \textbf{estabilización de la penalización} para $D \geq 8$, indicando que el sistema ha alcanzado la latencia de RAM pura.

La Tabla~\ref{tab:ram_comparativa} compara las tres variantes en RAM profunda ($L=3\,145\,728$, footprint 192 MB):

\begin{table}[H]
\centering
\caption{Comparativa en RAM profunda ($L=3\,145\,728$, footprint 192 MB)}
\label{tab:ram_comparativa}
\begin{tabular}{lcccl}
\toprule
\textbf{Stride} & \textbf{Double ind.} & \textbf{Int ind.} & \textbf{Double dir.} & \textbf{Mejor variante} \\
\midrule
D=1   & 7.22  & 6.81  & 7.16  & Int ind. ($-5\%$) \\
D=8   & 11.10 & 7.89  & 8.56  & Int ind. ($-\mathbf{29\%}$) \\
D=16  & 18.51 & 9.40  & 16.64 & Int ind. ($-\mathbf{49\%}$) \\
D=64  & 20.86 & 17.53 & 18.72 & Int ind. ($-16\%$) \\
D=128 & 19.68 & 19.71 & 16.59 & Double dir. ($-16\%$) \\
\bottomrule
\end{tabular}
\end{table}

La excepción notable es D=128 en RAM, donde \texttt{int} indirecto (19.71 ciclos) iguala al \texttt{double} indirecto. La razón es que con $R_\text{int} = 2 \times R_\text{double}$, el vector \texttt{ind[]} tiene el doble de elementos y ocupa proporciones masivas de memoria, generando accesos adicionales a RAM para leer los propios índices.

\textbf{Observación clave:} D=1 permanece completamente inmune a la jerarquía de memoria. Con 192 MB de datos en RAM, el resultado es $\sim$7.2 ciclos, prácticamente idéntico a la zona L1. Los cuatro prefetchers hardware actúan coordinadamente ocultando por completo la latencia de RAM.

% ════════════════════════════════════════════════════════════════════════════
\section{La ``escalera'' completa de la jerarquía}

Con los valores de L ampliados hasta 192 MB es posible visualizar la escalera completa de latencias. Para $D=64$ (double indirecto), los valores caracterizan cada nivel:

\begin{table}[H]
\centering
\caption{Escalera de latencias observadas para D=64, double indirecto}
\label{tab:escalera}
\begin{tabular}{llll}
\toprule
\textbf{Zona} & \textbf{L repr.} & \textbf{Footprint} & \textbf{Ciclos/acceso} \\
\midrule
L1          & 384       & 24 KB   & 7.15 \\
L2          & 15\,360   & 960 KB  & 7.41 (+0.3) \\
L3 entrada  & 40\,960   & 2.5 MB  & 8.18 (+0.8) \\
L3 profunda & 786\,432  & 48 MB   & 18.43 (+10.2) \\
RAM         & 3\,145\,728 & 192 MB & 20.86 (+2.4) \\
\bottomrule
\end{tabular}
\end{table}

La diferencia L1$\to$RAM real para D=64 es de $\Delta = 20.86 - 7.15 = \mathbf{13.7}$ ciclos, frente a los $\sim$145 ciclos teóricos de latencia pura de RAM. El prefetcher atenúa la penalización en un factor de aproximadamente $\mathbf{10\times}$.

La escalera clásica aparece, pero muy atenuada. El salto más pronunciado ocurre en la transición L2$\to$L3 dentro de la zona L3 ($\Delta \approx 10.2$ ciclos entre entrada y fondo de L3), mientras que la transición L3$\to$RAM añade solo $\sim$2.4 ciclos adicionales porque el prefetcher sigue actuando parcialmente incluso en RAM.

% ════════════════════════════════════════════════════════════════════════════
\section{Análisis comparativo de variantes}

\subsection{Double vs. Int (acceso indirecto)}

La hipótesis inicial era que ambas variantes deberían tener latencias similares para los mismos valores de L, porque el número de líneas referenciadas es idéntico. Los resultados confirman esto parcialmente, con diferencias importantes en determinados regímenes:

\begin{itemize}[leftmargin=1.5em]
  \item \textbf{Zona L1/L2:} \texttt{int} es $\sim$0.3 ciclos más rápido por la menor latencia de \texttt{ADD} vs. \texttt{FADD}. El número de líneas accedidas es el mismo.
  \item \textbf{Zona L3 media (D=16):} \texttt{int} es notablemente mejor ($\sim$7.3 vs. $\sim$10.3 ciclos para $L=163\,840$). Con \texttt{int}, $R$ es el doble pero el vector \texttt{ind[]} pesa la mitad relativa al footprint de \texttt{A[]}, reduciendo la contención.
  \item \textbf{RAM con D=16:} \texttt{int} obtiene 9.40 ciclos frente a 18.51 del \texttt{double}, una mejora del \textbf{49\%}. El IP Stride detecta mejor el patrón porque hay más elementos útiles por línea traída.
  \item \textbf{RAM con D=64/128:} las diferencias se reducen porque $R_\text{int}$ tan grande hace que \texttt{ind[]} sea enorme ($\sim$750 MB para $L=3\,145\,728$, $D=64$), añadiendo tráfico masivo a RAM solo para leer los índices.
\end{itemize}

\subsection{Directo vs. Indirecto (tipo Double)}

La eliminación del vector \texttt{ind[]} tiene dos efectos combinados: (1) se elimina un acceso de memoria por iteración (la lectura del índice antes del acceso a \texttt{A[]}), y (2) el espacio de caché liberado por \texttt{ind[]} queda disponible exclusivamente para \texttt{A[]}.

\begin{itemize}[leftmargin=1.5em]
  \item \textbf{Zona L1/L2:} La mejora con strides grandes (D=64: 6.42 vs. 7.15 ciclos en L1) se debe a que sin \texttt{ind[]}, L1 tiene más espacio para \texttt{A[]}.
  \item \textbf{Zona L3 media:} El acceso directo es mejor (7.57 vs. 10.32 ciclos para D=16, $L=163\,840$). La dependencia de datos \texttt{A[ind[i]]} serializa los accesos e impide que el prefetcher gestione dos streams de forma eficiente.
  \item \textbf{Zona RAM (D=8):} La mayor diferencia absoluta: 8.56 vs. 11.10 ciclos, mejora del \textbf{23\%}. El patrón directo \texttt{A[i*D]} es un stride perfectamente uniforme que el prefetcher detecta con máxima eficiencia.
\end{itemize}

% ════════════════════════════════════════════════════════════════════════════
\section{Resumen de los valores de latencia identificados}

\begin{table}[H]
\centering
\caption{Latencias observadas experimentalmente vs. latencias teóricas del procesador}
\label{tab:latencias}
\begin{tabular}{lrrl}
\toprule
\textbf{Nivel} & \textbf{Latencia teórica} & \textbf{Ciclos observados} & \textbf{Factor prefetcher} \\
\midrule
L1d  & $\sim$5 ciclos  & $\sim$7 ciclos  & — (incluye overhead del bucle) \\
L2   & $\sim$14 ciclos & $\sim$7 ciclos  & $2\times$ (oculta completamente) \\
L3   & $\sim$45 ciclos & 9--18 ciclos    & $2.5\times$--$5\times$ \\
RAM  & $\sim$150 ciclos & 11--21 ciclos   & $7\times$--$14\times$ \\
\bottomrule
\end{tabular}
\end{table}

Los valores observados en L1 ($\sim$7 ciclos) son superiores a la latencia teórica de 5 ciclos porque la compilación sin optimización (\texttt{-O0}) añade overhead en el acceso al array de índices y en el control del bucle.

% ════════════════════════════════════════════════════════════════════════════
\section{Conclusiones}

\begin{enumerate}[leftmargin=1.5em]
  \item \textbf{La experimentación cubre toda la jerarquía.} Con 11 valores de L desde 24 KB hasta 192 MB se caracterizan experimentalmente los cuatro niveles de la jerarquía: L1, L2, L3 y RAM. Los límites $S_1=768$, $S_2=20\,480$ y $S_3=786\,432$ líneas se validan directamente por los cambios de pendiente en las curvas.

  \item \textbf{D=1 es inmune a la jerarquía de memoria.} Con acceso completamente secuencial, los cuatro prefetchers hardware actúan coordinadamente ocultando las latencias de L2, L3 y RAM. Con 192 MB de datos, el coste es $\sim$7.2 ciclos, indistinguible del valor en L1.

  \item \textbf{La "escalera" clásica aparece pero muy atenuada.} La diferencia L1$\to$RAM para D=64 es de solo 13.7 ciclos en lugar de los $\sim$145 teóricos, ya que el prefetcher atenúa la penalización aproximadamente $10\times$.

  \item \textbf{La transición L3$\to$RAM es claramente observable} entre $L=786\,432$ y $L=1\,572\,864$ para todos los strides excepto D=1, confirmando experimentalmente que $S_3=786\,432$ líneas es el límite correcto de L3.

  \item \textbf{D=16 es el peor stride} en las variantes con acceso indirecto. El prefetch pollution del Spatial Prefetcher genera tráfico inútil que degrada el rendimiento hasta $\sim$19 ciclos en RAM, un factor $2.6\times$ respecto a D=1.

  \item \textbf{El acceso directo supera al indirecto} principalmente en L3 y RAM, con mejoras de hasta el 23\% en RAM para D=8. La eliminación de la dependencia de datos \texttt{ind[i]$\to$A[]} permite al prefetcher operar sobre un único stream regular en lugar de dos.

  \item \textbf{El tipo \texttt{int} supera al \texttt{double}} para strides pequeños y medios incluso en RAM. Para D=16 en RAM profunda obtiene 9.40 ciclos frente a 18.51 del \texttt{double} indirecto (mejora del 49\%). La razón es el menor peso relativo de \texttt{ind[]} y la mayor densidad de elementos útiles por línea de caché traída.

  \item \textbf{Los valores se estabilizan en RAM profunda.} Entre $L=1\,572\,864$ y $L=3\,145\,728$ los ciclos/acceso varían menos del 2\% en todas las variantes, confirmando que el sistema ha alcanzado la latencia de RAM pura sin degradación adicional.
\end{enumerate}

% ════════════════════════════════════════════════════════════════════════════
\appendix
\section{Fragmentos de código de los experimentos adicionales}

\subsection{Acceso directo: \texttt{acp1\_directo.c}}

\begin{lstlisting}[caption={Bucle de medición con acceso directo (sin vector de índices)}, label=lst:directo]
/* Sin ind[]: acceso directo A[i * D] */
start_counter();
for (int k = 0; k < REPS; k++) {
    double suma = 0.0;
    for (long long i = 0; i < R; i++)
        suma += A[i * D];   // Un único acceso a memoria por iteración
    acum[k] = suma;
}
double ciclos_totales = get_counter();
\end{lstlisting}

\subsection{Tipo Int: \texttt{acp1\_int.c}}

\begin{lstlisting}[caption={Diferencia clave en acp1\_int.c: tipo int y cálculo de R}, label=lst:int]
/* Con int (4 bytes), R es el DOBLE que con double (8 bytes)
   para el mismo número de líneas L                         */
long long R = (long long)L * CLS / (D * sizeof(int));

int *A   = (int*)aligned_alloc(CLS, N * sizeof(int));
int *ind = (int*)malloc(R * sizeof(int));
int  acum[REPS];

/* Bucle de medición: reducción de suma de enteros */
for (int k = 0; k < REPS; k++) {
    int suma = 0;
    for (long long i = 0; i < R; i++)
        suma += A[ind[i]];
    acum[k] = suma;
}
\end{lstlisting}

\end{document}